\documentclass[11pt]{beamer}
\usepackage[utf8]{inputenc}
\usepackage{amsmath, amssymb, graphicx, listings, xcolor}

% Presentation Theme
\usetheme{Madrid}
\usecolortheme{default}

% Code Block Styling
\lstset{
    language=Python,
    basicstyle=\ttfamily\tiny, % Tiny font to fit the code on the slide
    keywordstyle=\color{blue}\bfseries,
    commentstyle=\color{green!50!black},
    stringstyle=\color{orange},
    frame=single,
    backgroundcolor=\color{gray!5},
    showstringspaces=false
}

\title[Numerical Integration]{Astrodynamics 2026-1}
\subtitle{Algorithm Design: \\ Numerical Integration of Low-Thrust Trajectories}
\author{Prof. Juan}
\institute{Universidad de Antioquia - Aerospace Engineering}
\date{Spring 2026}

\begin{document}

% --- Title Slide ---
\begin{frame}
    \titlepage
\end{frame}

% ==========================================
% SECTION 1: THE STATE VECTOR
% ==========================================
\section{The 7D State Vector}

\begin{frame}{Expanding the State Vector}
    In Keplerian (impulsive) mechanics, the state vector has 6 dimensions (Position and Velocity). 
    
    \vspace{0.3cm}
    Because continuous thrust consumes propellant over time, mass becomes a dynamic variable. We must expand the state vector ($\boldsymbol{Y}$) to \textbf{7 dimensions}:

    \begin{equation}
        \boldsymbol{Y} = 
        \begin{bmatrix}
            x \\ y \\ z \\ v_x \\ v_y \\ v_z \\ m
        \end{bmatrix}
    \end{equation}
    
    \vspace{0.3cm}
    \textit{Mass is no longer a static constant; it is a coordinate that moves through time just like $x$ or $y$.}
\end{frame}

% ==========================================
% SECTION 2: THE EQUATIONS OF MOTION
% ==========================================
\section{The ODE Physics}

\begin{frame}{Constructing the Derivative (Physics)}
    To numerically integrate the trajectory, we must calculate the rate of change of the state vector: $d\boldsymbol{Y}/dt$.
    
    \vspace{0.3cm}
    \textbf{1. Gravity:}
    \begin{equation}
        \boldsymbol{a}_{grav} = -\frac{\mu}{|\boldsymbol{r}|^3} \boldsymbol{r}
    \end{equation}
    
    \textbf{2. Thrust (Tangential Steering):}
    \begin{equation}
        \boldsymbol{a}_{thrust} = \frac{T}{m(t)} \left( \frac{\boldsymbol{v}}{|\boldsymbol{v}|} \right)
    \end{equation}
    
    \textbf{3. Mass Depletion (Rocket Equation):}
    \begin{equation}
        \dot{m} = -\frac{T}{I_{sp} g_0}
    \end{equation}
\end{frame}

\begin{frame}{Constructing the Derivative (Output)}
    By combining these rates of change, our Ordinary Differential Equation (ODE) function must return the following 7D array at every time step $t$:
    
    \vspace{0.3cm}
    \begin{equation}
        \frac{d\boldsymbol{Y}}{dt} = 
        \begin{bmatrix}
            v_x \\ 
            v_y \\ 
            v_z \\ 
            a_{gx} + a_{tx} \\ 
            a_{gy} + a_{ty} \\ 
            a_{gz} + a_{tz} \\ 
            \dot{m}
        \end{bmatrix}
    \end{equation}
    
    \vspace{0.3cm}
    This is the exact mathematical engine operating inside GMAT's \texttt{FiniteBurn} routines!
\end{frame}

% ==========================================
% SECTION 3: PYTHON IMPLEMENTATION
% ==========================================
\section{Python Code}

\begin{frame}[fragile]{Python Implementation}
    We can write this exact ODE in Python and use SciPy's Runge-Kutta solver.
    
\begin{lstlisting}[language=Python]
import numpy as np
from scipy.integrate import solve_ivp

def low_thrust_ode(t, state, mu, thrust, isp, g0):
    # 1. Unpack state
    r_vec = state[0:3]
    v_vec = state[3:6]
    m = state[6]
    
    # 2. Gravity Acceleration
    r_mag = np.linalg.norm(r_vec)
    a_grav = (-mu / r_mag**3) * r_vec
    
    # 3. Thrust Acceleration (Tangential Steering)
    v_mag = np.linalg.norm(v_vec)
    u_v = v_vec / v_mag          # Unit velocity vector
    a_thrust = (thrust / m) * u_v
    
    # 4. Total Acceleration and Mass Depletion
    a_total = a_grav + a_thrust
    mdot = -thrust / (isp * g0)
    
    # 5. Return derivatives [dx/dt, dy/dt, dz/dt, dvx/dt, dvy/dt, dvz/dt, dm/dt]
    return np.concatenate((v_vec, a_total, [mdot]))

# --- Execution ---
# result = solve_ivp(low_thrust_ode, [0, max_t], initial_state, args=(mu, T, Isp, g0))
\end{lstlisting}
\end{frame}

% ==========================================
% SECTION 4: VALIDATION
% ==========================================
\section{Validation against GMAT}

\begin{frame}{Validation Exercise}
    \textbf{Your Assignment:}
    \begin{enumerate}
        \item Define a starting LEO orbit and a 0.5 N ion engine in Python.
        \item Integrate the trajectory for 30 days using \texttt{scipy.integrate.solve\_ivp}.
        \item Build the exact same spacecraft, \texttt{ElectricThruster}, and \texttt{ElectricTank} in GMAT.
        \item Propagate a \texttt{FiniteBurn} for 30 days.
    \end{enumerate}
    
    \vspace{0.5cm}
    Compare the final Semi-Major Axis ($a$) and final Spacecraft Mass ($m$). If your Python ODE is written correctly, it will perfectly match GMAT's numerical output!
\end{frame}

\end{document}